\documentclass[article,paper=b6]{jlreq}
\setlength{\columnseprule}{0.5pt}
\setlength{\columnsep}{5\zw}
\usepackage[left=20mm,right=20mm]{geometry}
\usepackage{amssymb,mathtools,mathcomp}
\usepackage{enumitem,tasks}
\usepackage{physics,physics2}
\settasks{label-width=2\zw}
\usephysicsmodule{ab,ab.braket}
\usepackage{bm}
\usepackage{graphicx,tabularx,here}
\usepackage{tikz,tikz-3dplot}
\usetikzlibrary{math,angles,quotes,calc,intersections}
\usepackage{kyotsutest}
\usepackage{answerbox}
\usepackage[pdftitle={証明},pdfauthor={まさし},hidelinks]{hyperref}
\pagestyle{empty}

\begin{document}
\noindent\textbf{［証明］}

道の面積は，道の外側の円の面積から道の内側の円の面積を引いたものであるから，
\begin{align*}
S &{}={} {\pi}\ab(r + a)^{2} - {\pi}{r}^{2}\\
&{}={} {\pi}\ab\{\ab({r}^{2}  + 2ar + {a}^{2}) - {r}^{2}\}\\
&{}={} {\pi}\ab(2ar + {a}^{2})\\
&{}={} {\pi}a\ab(2r + a)
\end{align*}
である。

また，道の真ん中を通る円の半径は $\ab(r + \dfrac{a}{2}){\,}\textrm{m}$ であるから，
\begin{align*}
\ell &{}={} 2{\pi}\ab(r + \dfrac{a}{2})\\
&{}={} {\pi}\ab(2r + a)
\end{align*}
であり，
$$a\ell = {\pi}a\ab(2r + a)$$
となる。これは道の面積$S$に等しい。

したがって，$S = a\ell$ である。
\end{document}